\documentclass{workbook}

\newcommand{\copyrightdate}{2024}

\usepackage{ifxetex}
\usepackage[utf8]{inputenc}
\usepackage{hyperref}
\usepackage{hyperxmp} % Embed meta data into the PDF
\hypersetup{%
	hidelinks=true,
	linkcolor = {0 0 1},
	% Metadata to be embedded by hyperxmp
	pdftitle={Ordinary Differential Equations (\jobname)},
	pdfauthor={Bernardo Galv\~ao-Sousa, Jason Siefken},
	pdfauthortitle={Author},
	pdfcopyright={Copyright (C) \copyrightdate, Bernardo Galv\~ao-Sousa, Jason Siefken},
	pdfsubject={Ordinary Differential Equations textbook/workbook},
	pdfkeywords={ordinary differential equations, ODEs, modeling, vectors, mathematics, textbook},
	pdfurl={https://github.com/siefkenj/IBLODEs/},
	pdflicenseurl={https://creativecommons.org/licenses/by-sa/4.0/},
}

%%%
% import all needed packages and macros
%%%
\usepackage[yyyymmdd]{datetime}
\input{common/preamble.tex}

\usepackage{breqn}

%%%
% Set up the footers to have the correct copyright notices
%%%

\fancypagestyle{siefken}{%
	\rfoot{\footnotesize\it \copyright\,Bernardo Galv\~ao-Sousa \& Jason Siefken, \copyrightdate \ \makebox(30,5){\includegraphics[height=1.2em]{by-sa.pdf}}}
	\lfoot{}
	\renewcommand{\headrulewidth}{0pt}
}


%%
% Allow hiding of environments
%%
\usepackage{environ}% http://ctan.org/pkg/environ
\makeatletter
\newcommand{\voidenvironment}[1]{%
  \expandafter\providecommand\csname env@#1@save@env\endcsname{}%
  \expandafter\providecommand\csname env@#1@process\endcsname{}%
  \@ifundefined{#1}{}{\RenewEnviron{#1}{}}%
}
\makeatother
% allow pagebreaks that only display in `standard` mode
\newcommand{\displayonlynewpage}{\begin{displayonly}\newpage\end{displayonly}}
% allow pagebreaks that only display in `book` mode
\newcommand{\bookonlynewpage}{\begin{bookonly}\newpage\end{bookonly}}


%
% Set up the three render modes: standard, instructor, and solutions.
% These render with varying amounts of extra data (like solutions and notes)
%
\newtoggle{instructor}
\newtoggle{standard}
\newtoggle{solutions}
\newtoggle{book}
\newtoggle{slides}
\newcommand{\setinstructor}{
	\toggletrue{instructor}
	\togglefalse{standard}
	\togglefalse{solutions}
	\togglefalse{book}
	\togglefalse{slides}
}
\newcommand{\setstandard}{
	\togglefalse{instructor}
	\toggletrue{standard}
	\togglefalse{solutions}
	\togglefalse{book}
	\togglefalse{slides}
}
\newcommand{\setsolutions}{
	\togglefalse{instructor}
	\togglefalse{standard}
	\toggletrue{solutions}
	\togglefalse{book}
	\togglefalse{slides}
}
\newcommand{\setbook}{
	\togglefalse{instructor}
	\togglefalse{standard}
	\togglefalse{solutions}
	\toggletrue{book}
	\togglefalse{slides}
}
\newcommand{\setslides}{
	\togglefalse{instructor}
	\togglefalse{standard}
	\togglefalse{solutions}
	\togglefalse{book}
	\toggletrue{slides}
}


%
% Infer the document level from the \jobname
%
\usepackage{xstring}
\IfSubStr{\jobname}{\detokenize{book}}{\setbook}{
	\IfSubStr{\jobname}{\detokenize{solutions}}{\setsolutions}{
		\IfSubStr{\jobname}{\detokenize{instructor}}{\setinstructor}{
			\IfSubStr{\jobname}{\detokenize{slides}}{\setslides}{
				\setstandard
			}
		}
	}
}


\setbookoptions{
	twosided = false,
	inline solutions = false,
}


\NewColoredEnvironment{
	name = lesson,
	display name = Lesson,
	banner color = Plum,
	title color = Plum,
	banner on left = true,
	open right = false,
}
\NewColoredEnvironment{
	name = module,
	display name = Module,
	banner color = Turquoise,
	title color = Cerulean,
	definition color = Cerulean,
	theorem color = myorange,
}
\NewColoredEnvironment{
	name = appendix,
	display name = Appendix,
	banner color = LimeGreen,
	title color = LimeGreen!70!Green!80!black,
	definition color = Cerulean,
	theorem color = myorange,
}
\NewColoredEnvironment{
	name = indices,
	display name = Indices,
	banner color = Green,
	title color = Green,
}
\NewColoredEnvironment{
	name = tutorial,
	display name = Tutorial,
	banner color = Peach,
	title color = Peach!80!black,
	emphbox color = Peach,
	% We will print tutorial worksheets back-to-back to save space
	open right = false,
}




\loadgeometry{default}

%
% Hide the non-problem environments
%
\newcommand{\coversubtitle}{} % we override the subtitle in each mode, so make sure the command exists to override.
\iftoggle{instructor}{
	\voidenvironment{module}
	\voidenvironment{appendix}
	\voidenvironment{bookonly}
	\voidenvironment{displayonly}
	\renewcommand{\coversubtitle}{Instructor Guide}
}{}
\iftoggle{solutions}{
	\voidenvironment{module}
	\voidenvironment{appendix}
	\voidenvironment{bookonly}
	\voidenvironment{displayonly}
	\voidenvironment{lesson}
	\voidenvironment{notes}
	\renewcommand{\coversubtitle}{Solutions}
}{}
\iftoggle{standard}{
	\voidenvironment{module}
	\voidenvironment{appendix}
	\voidenvironment{bookonly}
	\voidenvironment{solution}
	\voidenvironment{annotation}
	\voidenvironment{lesson}
	\renewcommand{\coversubtitle}{MAT244 Notes}
	\loadgeometry{default}
}{}
\iftoggle{book}{
	\voidenvironment{displayonly}
	\voidenvironment{solution}
	\voidenvironment{annotation}
	\voidenvironment{lesson}
	\renewcommand{\coversubtitle}{{\hspace{-5pt}\begin{tabular}{l}MAT244 Workbook\\\small\today{} Edition\end{tabular}}}
	\setbookoptions{
		twosided = true,
		inline solutions = false,
	}
	\loadgeometry{book}
}{}
\iftoggle{slides}{
	\voidenvironment{module}
	\voidenvironment{appendix}
	\voidenvironment{bookonly}
	\voidenvironment{solution}
	\voidenvironment{annotation}
	\voidenvironment{lesson}
	\renewcommand{\coversubtitle}{MAT244 Slides}
	\loadgeometry{slides}
	\initSlides
}{}
%\voidenvironment{solution}
%\voidenvironment{annotation}
%\voidenvironment{lesson}
%%\voidenvironment{notes}
%%\voidenvironment{displayonly}

% Allow an index to be created
\makeindex[title=Index of Terms, columns=3]
\makeindex[name=definitions, title=Index of Definitions, columns=3]
\makeindex[name=symbols, title=Index of Symbols, columns=3]

\indexsetup{
	level=\Heading,
	noclearpage
}

\begin{document}
%%
%% Import definitions from definition.tex; all definitions can be restated multiple times
%%

\input{common/definitions.tex}

%%
%% End Definitions
%%




% Needed to get different PDF bookmarks from the TOC entries
\hypersetup{bookmarksdepth=3}

\pagestyle{empty}

\input{common/cover.tex}

\newpage

\begin{bookonly}
	\clearpage
	\hbox{}
	\newpage
	\input{modules/preface.tex}
	\section*{Contributors}
	\input{common/contributors.tex}
	\section*{Dedication}
	\begin{center}
		This book is dedicated to
		\href{https://www.gazettetimes.com/news/local/obituaries/dr-robert-main-burton/article_9c087f07-c005-515a-bb3f-2c9c6a6b7332.html}{\color{blue}Dr.~Bob Burton}---friend and mentor.

		\emph{\large ``Sometimes you have to walk the mystical path with practical feet.''}
	\end{center}
	\newpage
	\mbox{}
	{
		\pagestyle{empty}
		\setcounter{tocdepth}{1}
		\tableofcontents
		\thispagestyle{empty}
	}
	\newpage
	\mbox{}
	\newpage
\end{bookonly}

\setcounter{page}{1}
\pagestyle{siefken}


\addcontentsline{toc}{chapter}{Lessons}

\begin{module}\label{module1}
	\Title{Modeling}

	In this module you will learn
	\begin{itemize}
		\item ??
	\end{itemize}

	\Heading{Modeling}

Suppose you are observing some \emph{green} ants walking on the sidewalk.
In the first minute you record 10 ants. In the second minute you
record 20. In the third minute, you record 40 ants. This continues until
there are too many ants for you to count.

\begin{center}
	\begin{tabular}{c|c}
		Minute & \#Green Ants\\
		\hline
		1 & 10\\
		2 & 20\\
		3 & 40\\
		4 & 80\\
		$\vdots$ & $\vdots$
	\end{tabular}
\end{center}

Since you lost count of the ants, you decide to use mathematics to try and figure out
how many ants walked by on minutes $5$, $6$, \ldots. You notice the pattern that
\[
	\text{Green ants per minute $n$} = 2^{n-1}\cdot 10.
\]
Stupendous! Mathematics now predicts there were $160$ ants during minute $5$. But something
else catches your eye. Across the sidewalk are \emph{brown} ants. You count these
ants every minute.

\begin{center}
	\begin{tabular}{c|c}
		Minute & \#Brown Ants\\
		\hline
		1 & 3\\
		2 & 6\\
		3 & 12\\
		4 & 24\\
		$\vdots$ & $\vdots$
	\end{tabular}
\end{center}

The pattern is slightly different. This time, 
\[
	\text{Green ants per minute $n$} = 2^{n-1}\cdot 3.
\]

Your friend, who was watching you the whole time, looks confused. ``Why come up with two complicated equations
when you can describe both types of ant at once?'' they declare.

\begin{center}
	\begin{tabular}{c}
		$\text{\#Ants at minute $n$}\ =\ 2\cdot(\text{\#Ants at minute $n-1$})$\\
		$\text{\#Green ants at minute 1}=10$\\
		$\text{\#Brown ants at minute 1}=3$\\
	\end{tabular}
\end{center}

Your friend has a point. Their model is elegant, but \emph{your} model can predict how many ants pass by at minute $3.222$!
Though, your friend would probably complain that $46.654$ is not a number of ants\ldots.

\medskip

You and your friend have just come up with two different \emph{mathematical models} for the number of ants
that walk across the sidewalk. They happen to make similar predictions for each minute and each have their
strengths and weakenesses. In this course, we will be focused on a particular type of mathematical model---one
that uses \emph{differential equations} at its core.

\Heading{Types of Models}

\begin{definition}[Mathematical Model]
A \emph{mathematical model}\index{model} is a description of the world
	\begin{enumerate}
		\item created in the service of answering a question, and
		\item where the complexity of the world has been abstracted away to numbers, quantities, and their relationships\footnote{ Other
mathematical objects are also allowed.}.
	\end{enumerate}
\end{definition}

In the previous situation, the \emph{question} you were trying to answer was ``how many ants are there at a given minute?''.
We sidestepped difficult issues like, ``Is an ant that is missing three legs still an ant?'' by using the common-sense
convention that ``the number of ants is a whole number and one colored blob that moves under its own power corresponds to one ant''; thus,
we could use single numbers to represent our quantity of interest (the ants).

You and your friend already came up with two types of models.
\begin{itemize}
	\item An \textbf{analytic} model based on known functions.
	\item A \textbf{recursive} model where subsequent terms are based on previous terms and initial conditions.
\end{itemize}

If we define $A(n)$ to be the number of ants corssing the sidewalk at minute $n$, the \emph{analytic} model presented for green ants is
\[
	A(n)=2^{n-1}\cdot 10
\]
and the \emph{recursive} model presented is
\begin{align*}
	A(1) &= 10\\
	A(n) &= 2\cdot A(n-1).
\end{align*}

Each type of model has pros and cons. For example, the analytic model allows you to calculate the number of ants at any minute
with few button presses on a calculator, whereas the recursive model is more difficult to calculate but
makes it clear that the number of ants is doubling every minute.

Often times recursive models are easier to write down than analytic models\footnote{ In fact, in many real-world situations, an analytic model
doesn't exist}, but they maybe harder to analyze. A third type of model has similarities to both analytic and recursive models, and
brings the power of calculus to modeling.

\begin{itemize}
	\item A \textbf{differential-equations} model is a model based on a relationship between a function's derivative(s), its values, and initial conditions.
\end{itemize}

\emph{Differential-equations} models are useful because derivatives correspond to rates of change---and things in the world are always changing.
Let's try to come up with a differential equations model for the ants.

We'd like an equation relating $A(n)$, the number of ants at minute $n$, to $A'(n)$, the \emph{instantaneous rate of change} of the number of ants at minute $n$.
Making a table, we see

\begin{center}
	\begin{tabular}{c|c|c}
		Minute & \#Brown Ants & Change (from prev. minute)\\
		\hline
		1 & 3 & ?\\
		2 & 6 & 3\\
		3 & 12 & 6\\
		4 & 24 & 12\\
	\end{tabular}
\end{center}

or

\begin{center}
	\begin{tabular}{c|c|c}
		Minute & \#Brown Ants & Change (from next minute)\\
		\hline
		1 & 3 & 3\\
		2 & 6 & 6\\
		3 & 12 & 12\\
		4 & 24 & ?\\
	\end{tabular}
\end{center}

depending on whether we record the change from the previous minute or up to the subsequent minute. Neither table gives the \emph{instantaneous} rate of
change, but in both tables, the change is proportional to the number of ants. So, we can set up a model
\[
	A'(n) = kA(n)
\]
where $k$ is a constant of proportionality that we will try to determine later. We've just written down a \emph{differential equation} with an undetermined parameter, $k$.

\begin{definition}[Differential Equation]
	A \emph{differential equation}\index{differential equation} is an equation relating a function to one or more of its derivatives.
\end{definition}

We'd like to figure out what $k$ is. One way to do so is to solve the differential equation and find the values of $k$ so that our model
correctly predicts the data. This is called \emph{fitting} the model to data.

\begin{definition}[Fitting a Model]
	Given a model $M$ with parameters $k_1$, $k_2$, $\ldots$ and data $D$, \emph{fitting the model $M$ to the data $D$}
	is the process of finding values for the parameters $k_1$, $k_2$, $\ldots$ so that $M$ most accurately predicts the data $D$.
\end{definition}

Note that, in general, fitting a model to data doesn't necessarily produce \emph{unique} values for the unknown parameters, and a fitted model
(especially when the data comes from real-world observations) usually doesn't reproduce the data exactly. However, in the case of these ants, we
just might get lucky.

\Heading{Solving Differential Equations}

In general, \emph{there is no algorithm for solving differential equations}. Fortunately, it is easy to check whether any particular
function is a solution to a differential equation, since there \emph{is} an algorithm to differentiate functions\footnote{ More specifically, there is an algorithm
to differentiate the \emph{elementary} functions, those functions formed by compositions, sums, products, and quotients of polynomials, trig, exponentials, and logs.}.
Because of this, \emph{guess and check} will be our primary method for solving differential equations.

\begin{example}
Use educated guessing to solve $A'(n)=kA(n)$.

Since $A'\approx A$, we might start with a function that is equal to its own derivative. There is a famous one: $e^n$. Testing, we see
\[
	\frac{\d}{\d n}e^n=e^n=ke^n
\]
if $k=1$, but it doesn't work for other $k$'s. Trying $e^{kn}$ instead yields
\[
	\frac{\d}{\d n}e^{kn}=ke^{kn}
\]
which holds for all $k$. Thus $A(n)=e^{kn}$ is \emph{a} solution to $A'(n)=kA(n)$. However, there are other solutions, because
\[
	\frac{\d}{\d n}Ce^{kn}=C\left(ke^{kn}\right)=k\left(Ce^{kn}\right),
\]
and so for every $C$, the function $A(n)=Ce^{kn}$ is a solution to $A'(n)=kA(n)$.
\end{example}

By guessing-and-checking, we have found an infinite number of solutions to $A'(n)=kA(n)$.  It's now time to fit our
model to the data.

\begin{example}
	Find values of $C$ and $k$ so that $A(n)=Ce^{kn}$ best models brown ants.

	Taking two rows from our brown ants table, we see
	\begin{align*}
		A(1) &= Ce^{k} = 3\\
		A(2) & = Ce^{2k}=6.
	\end{align*}

	Since $e^k$ can never be zero, from the first equation we get $C=3/e^k$. Combining with the second equation we find
	\[
		Ce^{2k}=\frac{3}{e^k}e^{2k}=3e^k=6
	\]
	and so $e^k=2$. In other words $k=\ln 2$. Plugging this back in, we find $C=3/2$. Thus our fitted model is
	\[
		A(n)=\tfrac{3}{2}e^{n\ln 2}.
	\]
\end{example}

Upon inspection, we can see that $\tfrac{3}{2}e^{n\ln 2} = 3\cdot 2^{n-1}$, which is the analytic model
that was first guessed for brown ants.



%One of the beauties of the guess-and-check method is that it \emph{cannot} produce an incorrect answer,
%since you're always checking (though it may fail to produce answers, any answer that it \emph{does} produce
%will be correct). Thus, even if you did some mathematically-dubious steps to produce a guess, if it passes the checks,
%the final answer is valid.
%
%\begin{example}
%Use Leibniz notation, integral calculus, and guess-and-check to come up with solutions to $A'(n)=kA(n)$.
%
%In Leibniz notation, $A'(n)$ is written as $\displaystyle\frac{\d A}{\d n}$, and we may suppress the function variable and write $A$ in place of $A(n)$. Thus, our equation becomes
%\[
%	\frac{\d A}{\d n}=kA.
%\]
%Treating $\d A/\d n$ as an actual ratio, we can bring all $A$'s to the left and $n$'s to the right, giving
%\[
%	\frac{1}{A}\d A = k\d n.
%\]
%Integrating both sides gives
%\[
%	\int \frac{1}{A}\d A = \ln A + C_1 \qquad =\qquad  \int k\d n = kn + C_2,
%\]
%and so
%\[
%	\ln A + C_1 = kn + C_2 \qquad \iff \qquad \ln A = kn + (C_2-C_1) = kn + C_3.
%\]
%Solving for $A$ by exponentiating both sides gives
%\[
%	A = e^{\ln A} = e^{kn + C_3} = C_4 e^{kn}.
%\]
%Much of what we've done in the preceding calculation \emph{was not rigorous mathematics}, but that's okay, because
%it leads us to the guess: $A(n)=Ce^{kn}$. Finally, we can test our guess and verify
%\[
%	\frac{\d}{\d n}Ce^{kn}=C\left(ke^{kn}\right)=k\left(Ce^{kn}\right),
%\]
%and so $A(n)=Ce^{kn}$ is a solution for every $C$.\footnote{ Notice that in our ``guess work'', $C_4$ was actually restricted to be a positive
%number, but our solution works with $C$ being positive or negative.}
%
%\end{example}




	%\input{modules/module1-exercises.tex}
\end{module}

%
% Hours 1-6
%

\begin{slide}
	\question
	You are observing starfish that made their way to a previously uninhabited tide-pool.
	You'd like to predict the year-on-year population of these starfish.

	You start with a simple assumption
	\[
		\# \text{new children per year}\ \sim\ \text{size of current population}
	\]
	\begin{parts}
		\item Come up with a mathematical model for the number of star fish in a given year.
		Your model should
		\begin{itemize}
			\item Define any notation (variables and parameters) you use
			\item Include at least one formula/equation
			\item Explain how your formula/equation relates to the starting assumption
		\end{itemize}
	\end{parts}
\end{slide}

\begin{slide}
	\question
		Let

		\begin{tabular}{rl}
			(Birth Rate) & $K=1.1$ children per starfish per year \\
			(Initial Pop.) & $P_0=10$ star fish
		\end{tabular}

		and define the model \textbf{M$_1$} to be the model for starfish population with
		these parameters.
	\begin{parts}
		\item Simulate the total number of starfish per year using Excel.
	\end{parts}
\end{slide}

\begin{slide}
	\question
	Recall the model \textbf{M$_1$} (from the previous question). 

	Define the model \textbf{M$_1^*$} to be
	\[
		P(t) = P_0 e^{0.742 t}
	\]
	\begin{parts}
		\item Are \textbf{M$_1$} and \textbf{M$_1^*$} different models or the same?
		\item Which of \textbf{M$_1$} or \textbf{M$_1^*$} is better?
		\item List an advantage and a disadvantage for each of \textbf{M$_1$} and \textbf{M$_1^*$}.
	\end{parts}
\end{slide}

\begin{slide}
	\question
	In the model \textbf{M$_1$}, we assumed the starfish had $K$ children at one point during the year.

	\begin{parts}
		\item Create a model \textbf{M$_n$} where the starfish are assumed to have $K/n$ children $n$ times per year (at regular intervals).
		\item Simulate the models \textbf{M$_1$}, \textbf{M$_2$}, \textbf{M$_3$} in Excel. Which grows fastest?
		\item What happens to \textbf{M$_n$} as $n\to\infty$?
	\end{parts}

\end{slide}

\begin{slide}
	\question
	Exploring \textbf{M$_n$}

	We can rewrite the assumptions of \textbf{M$_n$} as follows:
	\begin{itemize}
		\item At time $t$ there are $P_n(t)$ starfish.
		\item $P_n(0)=10$
		\item During the time interval $(t, t+1/n)$ there will be (on average) $K/n$ new children per starfish.
	\end{itemize}

	\begin{parts}
		\item Write an expression for $P_n(t+1/n)$ in terms of $P_n(t)$.
		\item Write an expression for $\Delta P$, the change in population from time $t$ to $t+\Delta t$.
		\item Write an expression for $\frac{\Delta P}{\Delta t}$.
		\item Write down a \emph{differential equation} relating $P'(t)$ to $P(t)$ where $\displaystyle P(t)=\lim_{n\to\infty} P_n(t)$.
	\end{parts}
\end{slide}

\begin{slide}
	\question
	Define the model \textbf{M$_\infty$} by
	\begin{itemize}
		\item $P(0)=10$
		\item $P'(t) = kP(t)$
	\end{itemize}
	and recall the model \textbf{M$_1$} defined by
	\begin{itemize}
		\item $P_1(0)=10$
		\item $P_1(t+1) = KP(t)$ for $t\geq 0$ years and $K=1.1$.
	\end{itemize}


	\begin{parts}
		\item If $k=K=1.1$, does the model \textbf{M$_\infty$} produce the same population estimates as \textbf{M$_1$}?
		\item Suppose that \textbf{M$_1$} accurately predicts the population. Can you find a value of $k$ so that \textbf{M$_\infty$}
		accurately predicts the population?
		\item What are some advantages and disadvantages of the models \textbf{M$_1$} and \textbf{M$_\infty$}?
	\end{parts}
\end{slide}

\begin{slide}
	\question
	After more observations, scientists notice a seasonal effect on starfish. They propose a new model called \textbf{S}:
	\begin{itemize}
		\item $P(0)=10$
		\item $P'(t) = k\cdot P(t)\cdot |\sin(2\pi t)|$
	\end{itemize}
	
	\begin{parts}
		\item What can you tell about the population (without trying to compute it)?
		\item Assuming $k=1.1$, estimate the population after 10 years.
		\item Assuming $k=1.1$, estimate the population after 10.3 years.
	\end{parts}
\end{slide}

\begin{slide}
	\question
	Consider the following argument:
	\begin{quote}
		\color{blue}
		At $t=0$, the change in population $\approx P'(0)=0$, so
		\[
			P(1) \approx P(0)+P'(0)\cdot 1 = P(0)=10.
		\]
		At $t=1$, the change in population $\approx P'(1)=0$, so
		\[
			P(2) \approx P(2)+P'(2)\cdot 1 = P(0)=10.
		\]
		And so on.

		So, the population of starfish remains constant.
	\end{quote}
	
	\begin{parts}
		\item Do you believe this argument? Can it be improved?
	\end{parts}
\end{slide}

\begin{slide}
	\question
	(Simulating \textbf{M$_\infty$} with different $\Delta$s)
	
	\begin{center}
	\begin{tabular}{l|l|l|l}
		Time & Pop. ($\Delta=0.1$) & Time & Pop. ($\Delta=0.2$)\\
		\hline
		0.0&	10&		0.0&	10\\
		0.1&	11.1&		0.2&	12.2\\
		0.2&	12.321&		0.4&	14.884\\
		0.3&	13.67631&	0.6&	18.15848\\
		0.4&	15.1807041&	0.8&	22.1533456
	\end{tabular}
	\end{center}

	\begin{parts}
		\item Compare $\Delta=0.1$ and $\Delta=0.2$. Which approximation grows faster?
		\item Graph the population estimates for $\Delta=0.1$ and $\Delta=0.2$ on the same plot. What does the graph show?

		\vspace{1cm}
		\item What $\Delta$s give the largest estimate for the population at time $t$?
		\item Is there a limit as $\Delta\to 0$?
	\end{parts}
\end{slide}

\begin{slide}
	(Simulating \textbf{M$_\infty$} with different $\Delta$s)
	
	% https://www.desmos.com/calculator/x7qbgv6mkc
	\includegraphics[width=2.5in]{compare-deltas.png}

	\begin{parts}
		\item Compare $\Delta=0.1$ and $\Delta=0.2$. Which approximation grows faster?
		\item Graph the population estimates for $\Delta=0.1$ and $\Delta=0.2$ on the same plot. What does the graph show?

		\vspace{1cm}
		\item What $\Delta$s give the largest estimate for the population at time $t$?
		\item Is there a limit as $\Delta\to 0$?
	\end{parts}
\end{slide}

\begin{slide}
	\question
	\label{many-models}
	Consider the following models for starfish growth
	\begin{itemize}
		\item[\textbf{M}] \# new children per year $\sim$ current population
		\item[\textbf{N}] \# new children per year $\sim$ resources available per individual
		\item[\textbf{O}] \# new children per year $\sim$ current population times the fraction of total resources remaining
	\end{itemize}

	\begin{parts}
		\item Guess what the population vs. time curves look like for each model. \label{model-guess-part}
		\item Create a differential equation for each model.
		\item Simulate population vs. time curves for each model (but pick a common initial population).
	\end{parts}
\end{slide}

\begin{slide}
	\question
	Recall the models
	\begin{itemize}
		\item[\textbf{M}] \# new children per year $\sim$ current population
		\item[\textbf{N}] \# new children per year $\sim$ resources available per individual
		\item[\textbf{O}] \# new children per year $\sim$ current population times the fraction of total resources remaining
	\end{itemize}

	\begin{parts}
		\item Determine which population grows fastest in the short term and which grows fastest in the long term.
		\item Are some models more sensitive to your choice of $\Delta$ when simulating?
		\item Are your simulations for each model consistently underestimates? Overestimates?
		\item Compare your simulated results with your guesses from question \ref{many-models}.\ref{model-guess-part}.
		What did you guess correctly? Where were you off the mark?
	\end{parts}
\end{slide}

%
% Hours 7-9
%


\begin{slide}
	\question
	A simple model for population growth has the form
	\[
		P'(t) = kP(t)
	\]
	where $k$ is the \emph{birth rate}.

	\begin{parts}
		\item Create a better model for population that includes both births and deaths.
	\end{parts}
\end{slide}

\begin{slide}
	\question
	The \emph{Lotka-Volterra Predator-Prey} models two populations, $F$ (foxes) and $R$ (rabbits), simultaneously. It takes the form
	\begin{align*}
		F'(t) &= (B_F - D_F)\cdot F(t)\\
		R'(t) &= (B_R - D_R)\cdot R(t)
	\end{align*}
	where $B_{?}$ stands for births and $D_{?}$ stands for deaths.

	\begin{parts}
		\item Come up with appropriate formulas for $B_F$, $B_R$, $D_F$, and $D_R$, given your knowledge about
		how foxes and rabbits might interact.
		\item When would $B_F$ and $D_F$ be at their max/min?
		\item When would $B_R$ and $D_R$ be at their max/min?
	\end{parts}
\end{slide}

\begin{slide}
	\question
	Suppose the population of $F$ (foxes) and $R$ (rabbits) evolves over time following the rule
	\begin{align*}
		F'(t) &= (0.01\cdot R(t) - 1.1)\cdot F(t)\\
		R'(t) &= (1.1 - 0.1\cdot F(t))\cdot R(t)
	\end{align*}

	\begin{parts}
		\item Simulate the population of foxes and rabbits with a spreadsheet.
		\item Do the populations continue to grow/shrink forever? Are they cyclic?
		\item Should the humps/valleys in the rabbit and fox populations be in phase? Out of phase?
	\end{parts}
\end{slide}

\begin{slide}
	\question
	% https://utoronto-my.sharepoint.com/:x:/g/personal/jason_siefken_utoronto_ca/Eay4QOMvy7lNr5pOKRv22NgBLGUw7qMpSCShUjeAdrhsHQ?e=bpg4CP
	Open the spreadsheet

	\url{https://uoft.me/foxes-and-rabbits}

	which contains an Euler approximation for the Foxes and Rabbits population.
	\begin{align*}
		F'(t) &= (0.01\cdot R(t) - 1.1)\cdot F(t)\\
		R'(t) &= (1.1 - 0.1\cdot F(t))\cdot R(t)
	\end{align*}

	\begin{parts}
		\item Is the max population of the rabbits being over/under estimated? Sometimes over, sometimes under?
		\item What about the foxes?
		\item What about the min populations?
	\end{parts}
\end{slide}

\begin{slide}
	\question
	% https://utoronto-my.sharepoint.com/:x:/g/personal/jason_siefken_utoronto_ca/Eay4QOMvy7lNr5pOKRv22NgBLGUw7qMpSCShUjeAdrhsHQ?e=bpg4CP
	Open the spreadsheet

	\url{https://uoft.me/foxes-and-rabbits}

	which contains an Euler approximation for the Foxes and Rabbits population.
	\begin{align*}
		F'(t) &= (0.01\cdot R(t) - 1.1)\cdot F(t)\\
		R'(t) &= (1.1 - 0.1\cdot F(t))\cdot R(t)
	\end{align*}

	\begin{definition}[Component Graph \& Phase Plane]
		For a differential equation involving the functions $F_1$, $F_2$, \ldots, $F_n$, and the variable $t$,
		the \emph{component graphs} are the $n$ graphs of $(t, F_1(t))$, $(t, F_2(t))$, \ldots.
		
		The \emph{phase plane} or \emph{phase space} associated with the differential equation
		is the $n$-dimensional space with axes corresponding to
		the values of $F_1$, $F_2$, \ldots, $F_n$.
	\end{definition}

	\begin{parts}
		\item Plot the Fox vs.~Rabbit population in the \emph{phase plane}. 
		\item Should your plot show a closed curve or a spiral?
		\item What ``direction'' do points move along the curve as time increases? Justify by referring to the model.
		\item What is easier to see from plots in the phase plane than from component graphs (the graphs of
		fox and rabbit population vs. time)?
	\end{parts}
\end{slide}

\begin{slide}
	\question
	% https://utoronto-my.sharepoint.com/:x:/g/personal/jason_siefken_utoronto_ca/Eay4QOMvy7lNr5pOKRv22NgBLGUw7qMpSCShUjeAdrhsHQ?e=bpg4CP
	Open the spreadsheet

	\url{https://uoft.me/foxes-and-rabbits}

	which contains an Euler approximation for the Foxes and Rabbits population.
	\begin{align*}
		F'(t) &= (0.01\cdot R(t) - 1.1)\cdot F(t)\\
		R'(t) &= (1.1 - 0.1\cdot F(t))\cdot R(t)
	\end{align*}

	\begin{definition}[Equilibrium Solution]
	An \emph{equilibrium solution} to a differential equation or system of differential equations is a
	solution that is constant in the independent variable(s).
	\end{definition}

	\begin{parts}
		\item By changing initial conditions, what is the ``smallest'' curve you can get in the phase plane? What happens at
		those initial conditions?
		\item What should $F'$ and $R'$ be if $F$ and $R$ are \emph{equilibrium solutions}?
		\item How many equilibrium solutions are there for the fox-and-rabbit system? Justify your answer.
		\item What do the equilibrium solutions look like in the phase plane? What about their component graphs?
	\end{parts}
\end{slide}

\begin{slide}
	\question
	Recall the logistic model for starfish growth:
	\begin{itemize}
		\item[\textbf{O}] \# new children per year $\sim$ current population times the fraction of total resources remaining
	\end{itemize}
	which can be modeled with the equation
	\[
		P'(t) = k\cdot P(t)\cdot \left(1-\tfrac{R_i}{R}\cdot P(t)\right)
	\]
	where 
	\begin{itemize}
		\item $P(t)$ is the population at time $t$
		\item $k$ is a constant of proportionality 
		\item $R$ is the total number of resources
		\item $R_i$ is the resources that one starfish wants to consume
	\end{itemize}

	Use $k=1.1$, $R=1$, and $R_i=0.1$ unless instructed otherwise.

	\begin{parts}
		\item What are the equilibrium solutions for model \textbf{O}?
		\item What does a ``phase plane'' for model \textbf{O} look like? What do graphs of equilibrium solutions look like?
		\item Classify the behaviour of solutions that lie \emph{between} the equilibrium solutions. E.g., are they increasing, decreasing, oscillating?
	\end{parts}
\end{slide}

\begin{slide}
	\question

	\begin{definition}[Classification of Equilibria]
		An equilibrium solution \emph{$f$} is called
		\begin{itemize}
			\item \emph{attracting} if solutions locally converge to $f$
			\item \emph{repelling} if solutions locally diverge from $f$
			\item \emph{stable} if solutions do not locally diverge from $f$
			\item \emph{unstable} if solutions do not locally converge to $f$
			\item \emph{semi-stable} if solutions locally converge to $f$ from one side and
			locally diverge from $f$ on another.
		\end{itemize}
	\end{definition}

	Let
	\[
		F'(t) =\ ?
	\]
	be an unknown differential equation with equilibrium solution $f(t)=1$.

	\begin{parts}
		\item Draw an example of what solutions might look like if $f$ is \emph{attracting}.
		\item Draw an example of what solutions might look like if $f$ is \emph{repelling}.
		\item Draw an example of what solutions might look like if $f$ is \emph{stable}.
		\item Could $f$ be stable but \emph{not} attracting?
	\end{parts}
\end{slide}

\begin{slide}
	\question
	\begin{definition}[Classification of Equilibria]
		An equilibrium solution \emph{$f$} is called
		\begin{itemize}
			\item \emph{attracting} if solutions locally converge to $f$
			\item \emph{repelling} if solutions locally diverge from $f$
			\item \emph{stable} if solutions do not locally diverge from $f$
			\item \emph{unstable} if solutions do not locally converge to $f$
			\item \emph{semi-stable} if solutions locally converge to $f$ from one side and
			locally diverge from $f$ on another.
		\end{itemize}
	\end{definition}
	
	Recall the starfish population model \textbf{O} given by
	\[
		P'(t) = k\cdot P(t)\cdot \left(1-\tfrac{R_i}{R}\cdot P(t)\right)
	\]
	Use $k=1.1$, $R=1$, and $R_i=0.1$ unless instructed otherwise.


	\begin{parts}
		\item Classify the equilibrium solutions for model \textbf{O} as attracting/repelling/stable/unstable/semi-stable.
		\item Does changing $k$ change the nature of the equilibrium solutions? How can you tell?
	\end{parts}
\end{slide}

%
% Hours 10-12
%

\begin{slide}
	% https://www.desmos.com/calculator/ghavqzqqjn

	\question
	
	\begin{center}
	\includegraphics[width=2.5in]{slope-field-for-O.png}
	\end{center}

	A \emph{slope field} is a plot of small segments of tangent lines
	to solutions of a differential equation at different initial conditions.
	
	On the left is a slope field for model \textbf{O}, available at

	\url{https://www.desmos.com/calculator/ghavqzqqjn}

	\begin{parts}
		\item If you were sketching the slope field for model \textbf{O} by hand, what line would you sketch
		(a segment of) at $(5,3)$? Write an equation for that line.
		\item How can you recognize equilibrium solutions in a slope field?
		\item Describe different solutions to the \emph{differential equation} using words. Do all
		of those solutions make sense in terms of \emph{model \textbf{O}}?
	\end{parts}
\end{slide}

\begin{slide}
	% https://www.desmos.com/3d/fsfbhvy2h9

	\question
	
	\begin{center}
		\includegraphics[width=2in]{slope-field-3d.png}
	\end{center}

	3d slope fields are possible, but hard to interpret.

	On the left is a slope field for the Foxes--Rabbits model.

	\url{https://www.desmos.com/3d/fsfbhvy2h9}

	\begin{parts}
		\item What are the three dimensions in the plot?
		\item What should the graph of an equilibrium solution look like?
		\item What should the graph of a typical solution look like?
		\item What are ways to simplify the picture so we can more easily analyze solutions?
	\end{parts}
\end{slide}


\begin{slide}
	% https://www.desmos.com/calculator/wdgtznxndp
	%
	% Without Euler's built in:
	% https://www.desmos.com/calculator/vrk0q4espx

	\question
	
	\begin{center}
		\includegraphics[width=2.5in]{phase-portrait-foxes-wolves.png}
	\end{center}


	\begin{definition}[Phase Portrait]
		A \emph{phase portrait} or \emph{phase diagram} is the plot of a vector field in phase space
		where each vector rooted at $(x,y)$ is tangent to a solution curve passing through $(x,y)$
		and its length is given by the speed of a solution passing through $(x,y)$.
	\end{definition}

	On the left is a phase portrait for the Foxes--Rabbits model.

	\url{https://www.desmos.com/calculator/vrk0q4espx}

	\begin{parts}
		\item What do the $x$ and $y$ axes correspond to?
		\item Identify the equilibria in the phase portrait. What are the lengths of the vectors at those points?
		\item Classify each equilibrium as stable/unstable.
		\item Why is the vector at $(5,100)$ longer than the vector at $(10,100)$? Justify numerically.
	\end{parts}
\end{slide}

\begin{slide}
	% https://www.desmos.com/calculator/wdgtznxndp
	%
	% Without Euler's built in:
	% https://www.desmos.com/calculator/vrk0q4espx

	\question
	Sketch your own vector field where the corresponding system of differential equations:

	\begin{parts}
		\item Has an attracting equilibrium solution.
		\item Has a repelling equilibrium solution.
		\item Has no equilibrium solutions.
	\end{parts}
\end{slide}

\begin{slide}
	% https://www.desmos.com/calculator/ghavqzqqjn

	\question
	
	\begin{center}
	\includegraphics[width=2.5in]{slope-field-for-O.png}
	\end{center}

	Recall the slope field for model \textbf{O}.

	\begin{parts}
		\item What would a phase portrait for model \textbf{O} look like? Draw it.
		\item Where are the arrows the longest? Shortest?
		\item How could you tell from a 1d phase portrait whether an equilibrium solution is
		attracting/repelling/etc.?
	\end{parts}
\end{slide}

\begin{slide}
	%
	% Completed phase portrait in desmos:
	% https://www.desmos.com/calculator/tvjag852ja
	\question
	\label{q-phase}
	The following differential equation models the life cycle of a tree.
	In the model
	\begin{itemize}
		\item $H(t)\ =$ height (in meters) of tree trunk at time $t$
		\item $A(t)\ =$ surface area (in square meters) of all leaves at time $t$
	\end{itemize}
	\begin{align*}
		H'(t) &= 0.3\cdot A(t)-b\cdot H(t)\\
		A'(t) &= -0.3\cdot (H(t))^2 + A(t)
	\end{align*}
	and $0 \leq b \leq 2$

	\bigskip
	\phantom{x}
	\begin{parts}
		\item 
		\label{part-phase}
		Modify 

		{\small
		\url{https://www.desmos.com/calculator/vrk0q4espx}
		}

		to make a phase portrait for the tree model.
		\item What do equilibrium solutions mean in terms of tree growth?
		\item For $b=1$ what are the equilibrium solution(s)?
	\end{parts}
	\bigskip
	\phantom{x}
\end{slide}

\begin{slide}
	%
	% Completed phase portrait in desmos:
	% https://www.desmos.com/calculator/tvjag852ja
	\question
	The following differential equation models the life cycle of a tree.
	In the model
	\begin{itemize}
		\item $H(t)\ =$ height (in meters) of tree trunk at time $t$
		\item $A(t)\ =$ surface area (in square meters) of all leaves at time $t$
	\end{itemize}
	\begin{align*}
		H'(t) &= 0.3\cdot A(t)-b\cdot H(t)\\
		A'(t) &= -0.3\cdot (H(t))^2 + A(t)
	\end{align*}
	and $0 \leq b \leq 2$

	\bigskip
	\phantom{x}
	\begin{parts}
		\item Fix a value of $b$ and use a spreadsheet to simulate some solutions with different initial conditions.
		Plot the results on your phase portrait from \ref{q-phase}.\ref{part-phase}.
		\item What will happen to a tree with $(H(0), A(0))=(20,10)$? Does this depend on $b$?
		\item What will happen to a tree with $(H(0), A(0))=(10,10)$? Does this depend on $b$?
	\end{parts}
	\bigskip
	\phantom{x}
\end{slide}

\begin{slide}
	\question
	The tree model
	\begin{align*}
		H'(t) &= 0.3\cdot A(t)-b\cdot H(t)\\
		A'(t) &= -0.3\cdot (H(t))^2 + A(t)
	\end{align*}
	was based on the premises
	\begin{itemize}[leftmargin=3em]
		\item[ $P_{\text{height 1}}$] CO$_2$ is absorbed by the leaves and turned directly into trunk height.
		\item[ $P_{\text{height 2}}$] The tree is in a swamp and constantly sinks at a speed proportional to its height.
		\item[ $P_{\text{leaves 1}}$] Leaves grow proportionality to the energy available.
		\item[ $P_{\text{energy 1}}$] The tree absorbs energy from the sun proportionality to the leaf area.
		\item[ $P_{\text{energy 2}}$] It costs energy proportional to the square of the height for the tree to maintain its current size.
	\end{itemize}

	\begin{parts}
		\item How are the premises expressed in the differential equations?
		\item What does the parameter $b$ represent?
		\item Applying Euler's method to this system shows solutions that pass from the 1st to 4th quadrants of the phase plane.
			Is this realistic? Describe the life cycle of such a tree?
	\end{parts}
\end{slide}

\begin{slide}
	\question
	Recall the tree model
	\begin{align*}
		H'(t) &= 0.3\cdot A(t)-b\cdot H(t)\\
		A'(t) &= -0.3\cdot (H(t))^2 + A(t)
	\end{align*}

	\begin{parts}
		\item Find all equilibrium solutions for $0\leq b\leq 2$.
		\item For which $b$ does a tree have the possibility of living forever? If the wind occasionally blew off a few random leaves,
		would that change your answer?
		\item
		Find a value $b_5$ of $b$ so that there is an equilibrium with $H=5$.

		Find a value $b_{12}$ of $b$ so that there is an equilibrium with $H=12$.
		
		\item
		Predict what happens to a tree near equilibrium in condition $b_5$ and a tree near equilibrium in condition $b_{12}$.

	\end{parts}
\end{slide}

%
% Hours 13-14
%

\begin{slide}
	\question
	Consider the system of differential equations
	\begin{align*}
		x'(t) &= x(t)\\
		y'(t) &= 2y(t)
	\end{align*}

	\begin{parts}
		\item Make a phase portrait for the system.
		\item What are the equilibrium solution(s) of the system?
		\item 
		Find a formula for $x(t)$ and $y(t)$ that satisfy the initial conditions $(x(0), y(0))=(x_0, y_0)$.
		\item Let $\vec r(t)=(x(t),y(t))$. Find a matrix $A$ so that the differential equation can be equivalently expressed
		as
		\[
			\vec r'(t) = A\,\vec r(t).
		\]
		\item Write a solution to $\vec r' = A\,\vec r$ (where $A$ is the matrix you came up with).
	\end{parts}
\end{slide}








\begin{bookonly}
\begin{appendix}\label{APPSLEI}
	\Title{Systems of Linear Equations I}

	In this appendix you will learn
	\begin{itemize}
		\item What a system of linear equations is.
		\item What the solution set to a system of equations is, and what it means for a system of equations
			to be consistent or inconsistent.
		\item How augmented matrices can be used to solve systems of linear equations.
		\item How to apply row reduction to find a unique solution to a system of linear
			equations and to determine if a system of linear equations is consistent or inconsistent.
	\end{itemize}

	\input{modules/appendix1.tex}
	\input{modules/appendix1-exercises.tex}
\end{appendix}


\begin{appendix}
	\PrintExerciseSolutions
\end{appendix}
	\begin{indices}*
		\Title{Indices}

		\printindex[symbols]

		\bigskip
		\printindex

		\bigskip
		\printindex[definitions]
	\end{indices}
\end{bookonly}

\end{document}
